\documentclass[10pt,a4paper]{article}
\usepackage[margin=2cm]{geometry}
\usepackage[T1]{fontenc}
\usepackage{lmodern,booktabs,tabularx}
\usepackage[hidelinks]{hyperref}
\usepackage{microtype}
\setlength{\parindent}{0pt}
\setlength{\parskip}{5pt}
\title{GRACE installation and A100 training validation on UMT}
\author{MLIP-nucleation project}
\date{17 October 2026}
\begin{document}
\maketitle

\section{Accepted environment}
A separate environment was installed in
\path{/home6/u3129/mlip-research/envs/grace_ft_a100}.
The accepted versions are Python 3.11.17, TensorPotential 0.6.1,
TensorFlow 2.20.0, tf-keras 2.20.1, NumPy 2.4.6, ASE 3.29.0 and dpdata 1.1.0.
\texttt{pip check} reported no broken requirements. Existing DeepMD environments
were not modified. The prefix is on compute-writable home6 storage; the main
home filesystem is read-only on tested compute nodes.

Inside a suitable Slurm allocation, source:
\begin{verbatim}
source /home6/u3129/mlip-research/envs/grace_ft_a100/activate-grace.sh
gracemaker input.yaml
\end{verbatim}
The helper activates the prefix, selects legacy Keras, enables GPU memory growth,
and sets \texttt{XLA\_FLAGS} to the installed CUDA wheel's data root. It also
prepends the wheel's \texttt{ptxas} directory to \texttt{PATH}.
The CUDA data root is \path{lib/python3.11/site-packages/nvidia/cuda_nvcc}
relative to the environment. Its \path{nvvm/libdevice/libdevice.10.bc} is
needed even for some ordinary TensorFlow GPU operations when explicit
training JIT compilation is disabled. The successful run records hashes of
both \texttt{libdevice} and \texttt{ptxas}.

\section{Successful GPU check}
Slurm job \textbf{20950274} completed with exit code 0 on
\texttt{tesla-a101}, partition \texttt{v100}, account \texttt{imet-2026-2},
QoS \texttt{normal}. It requested one A100, four CPUs, 16 GiB RAM and a
five-minute limit. Scheduler elapsed time was \textbf{38 seconds}.
The device was an NVIDIA A100-PCIE-40GB with driver 580.126.20.

A reduced-size \texttt{GRACE\_1LAYER\_latest} model with 7113 trainable
parameters completed \textbf{five epochs and 20 optimizer updates}.
The input used angular cutoff 1, radial width 8, embedding width 8,
batch size 2, Adam and a cosine learning-rate schedule. Training JIT was
disabled for this bounded check. Both training and validation metrics were
finite. The regular checkpoint contains step 20 and only finite floating-point
tensors. GRACE also wrote its automatic \texttt{final\_model} SavedModel.

The runtime probe executed matrix and square-root operations and their
gradients on GPU:0. Independently, nine NVIDIA process observations identified
the training child PID on the allocated GPU. Thus a visible device alone
was not accepted as evidence of GPU training.

Inputs are explicitly synthetic harmonic-pair labels for four Al-labelled
atoms: eight independently generated training configurations and four validation
configurations, seed 20261017. They are neither DFT labels nor foundation-model
predictions and do not select the pilot material. Energies are in eV and forces
in eV/angstrom. The optional stored ASE stress is $-W/V$, where
$W=\sum_i\mathbf r_i\otimes\mathbf F_i$, in eV/angstrom$^3$ and Voigt order
$xx,yy,zz,yz,xz,xy$. Stress was not fitted in this check.

\section{Attempts and diagnosed failures}
\begin{center}
\begin{tabularx}{\linewidth}{@{}llllX@{}}
\toprule
Job & Kind & CPUs/GPU & Elapsed & Result \\
\midrule
20950196 & Install 01 & 4/0 & 0:07 & Failed: Conda plugin cache targeted read-only main home \\
20950197 & Install 02 & 4/0 & 30:26 & Timeout during extraction after slow package downloads \\
20950266 & Install 03 & 1/0 & 3:33 & Passed: full reinstall using retained wheel cache \\
20950272 & GPU check 01 & 4/1 & 0:28 & Failed: TensorFlow could not discover packaged libdevice \\
20950273 & GPU check 02 & 4/1 & 0:21 & Failed: minimal input omitted required scheduler\_params \\
20950274 & GPU check 03 & 4/1 & 0:38 & Passed: explicit CUDA paths and scheduler configuration \\
\bottomrule
\end{tabularx}
\end{center}

All three GPU attempts consumed a combined \textbf{87 single-GPU seconds},
within the initial ten GPU-minute envelope. Installation initially had a
30-minute CPU limit; Slurm denied an attempted live increase. The continuation
used one CPU and a one-hour limit, and finished in 3 minutes 33 seconds.
It reinstalled from cache because a timeout during extraction can leave
partially written files even when some package metadata already exists.
Installer HOME and caches were isolated on home6; no global shell configuration
or existing scientific environment was changed.

\section{Locks, records and actual checkpoint backup}
\path{environments/grace-a100.lock.json} indexes the accepted package set,
the installation reports and hashes. The explicit Conda lock is
\path{environments/grace-a100-conda-linux-64.lock.txt}; exact HTTPS wheel
URLs and SHA-256 values for 82 pip artifacts are in
\path{environments/grace-a100-pip-linux-64.lock.txt}.
Rebuild a new Linux x86\_64 prefix from the Conda lock first, then install the
pip lock with \texttt{--no-deps --require-hashes}. The activation helper is
prefix-specific. These locks describe software reconstruction; they are not
a backup backend for future scientific datasets and foundation checkpoints.

Run records are in \path{runs/grace-install-20261017-01/} through
\path{runs/grace-install-20261017-03/} and
\path{runs/grace-smoke-20261017-01/} through
\path{runs/grace-smoke-20261017-03/}. They preserve job IDs, submission intents,
input hashes, code revisions, exact training inputs, synthetic structures,
package artifacts, scheduler outcomes, logs and metrics. Collected files are
verified against remote SHA-256 values. The accepted training inputs are
identified by \path{runs/grace-smoke-20261017-03/manifest.json}.

The accepted regular checkpoint is also backed up as actual binary files in
\path{runs/grace-smoke-20261017-03/checkpoint/}, alongside the model definition
and \path{checkpoint-backup.json}. Its index is 3859 bytes and its data file
is 96650 bytes. This approximately 100 kB software-test artifact is small
enough for Git. Full checkpoint hashes and the auto-export's hashes are in
the run result; exported SavedModel files remain on the cluster.
The test checkpoint is not a selected scientific foundation model.

\section{Validation limits and next steps}
This confirms a small single-GPU GRACE-1L fit from scratch, finite saved weights,
and a usable isolated environment. It does not validate physical accuracy,
production throughput or memory scaling, K40m execution, multi-GPU execution,
XLA-compiled training, restart execution or numerical equivalence of the
automatic export. No foundation checkpoint was downloaded or fine-tuned.
Select trainable weights separately, recording source revision, SHA-256,
license, chemical coverage and DFT reference conventions, then validate the
real fine-tuning and deployment path with grouped scientific data splits.

\section{Primary references}
Accessed 17 October 2026. Runtime conclusions above are supported by the saved
cluster records rather than inferred from documentation.
\begin{itemize}
\item GRACE installation and legacy Keras guidance:
\url{https://gracemaker.readthedocs.io/en/latest/gracemaker/install/}.
\item GRACE training input schema:
\url{https://gracemaker.readthedocs.io/en/latest/gracemaker/inputfile/}.
\item Published TensorPotential 0.6.1 package metadata:
\url{https://pypi.org/pypi/tensorpotential/0.6.1/json}.
\item TensorFlow issue documenting libdevice discovery:
\url{https://github.com/tensorflow/tensorflow/issues/56927}.
\end{itemize}
\end{document}
